\chapter{Conclusion}
\label{ch:conclusion}

\section{Bilan}

L'IA a ete utilisee sur \textbf{deux faces} :
\begin{itemize}
  \item \textbf{Developpement} : ChatGPT pour amorcer le cahier des prompts ;
        editeur \textbf{Cursor} avec modeles \textbf{Claude} et \textbf{Grok}
        pour implementer et corriger le code de facon iterative ;
  \item \textbf{Produit} : Ollama et APIs (HF, Sightengine, TextGears, Gemini)
        plus des modeles ML locaux pour des besoins metier precis.
\end{itemize}

\section{Limites}

\begin{itemize}
  \item les modeles peuvent proposer du code incorrect : verification et prompts
        de correction sont indispensables ;
  \item les APIs gratuites ont des quotas et peuvent etre indisponibles ;
  \item les captures et prompts doivent etre reels (pas reconstitues).
\end{itemize}

\section{Perspectives}

Enrichir le journal de prompts avec toutes les sessions (en indiquant Claude ou
Grok), completer les figures, et ajouter d'autres captures des fonctionnalites
IA en situation reelle d'utilisation.
